\section{September 11, 2026}

\subsection{The subspace criterion}

\begin{proof}[Proof of Proposition~\ref{prop:subspace-criterion}]
  If $U$ is a subspace of $V$, then $U$ is a vector space under the
  operations inherited from $V$. It therefore contains the additive
  identity and is closed under addition and scalar multiplication.

  Conversely, suppose conditions (i)--(iii) hold. The additive identity
  belongs to $U$ by (i). If $u\in U$, then condition (iii), applied with
  the scalar $-1$, gives
  \[
    -u=(-1)u\in U.
  \]
  Hence $U$ contains the additive inverse of each of its elements. The
  remaining vector-space axioms are identities already valid in $V$, so
  they remain valid when the vectors are restricted to $U$. Thus $U$ is
  a subspace of $V$.

  Finally, if $U$ is nonempty and satisfies the two closure conditions,
  choose $u\in U$. Closure under scalar multiplication gives
  $0=0u\in U$, so condition (i) follows automatically.
\end{proof}

\subsection{Kernels and images}

\begin{definition}[Kernel and image]
  Let $T\in\mathcal L(V,W)$. The \emph{kernel} and \emph{image} of $T$
  are
  \[
    \ker T=\{v\in V:T(v)=0\}
  \]
  and
  \[
    \operatorname{im}T=\{T(v):v\in V\},
  \]
  respectively. The image is also called the \emph{range} of $T$.
\end{definition}

\begin{proposition}
  If $T\in\mathcal L(V,W)$, then $\ker T$ is a subspace of $V$ and
  $\operatorname{im}T$ is a subspace of $W$.
\end{proposition}

\begin{proof}
  Since $T(0)=0$, both sets contain the appropriate zero vector. If
  $u,v\in\ker T$ and $\alpha\in\bb{F}$, then
  \[
    T(u+v)=T(u)+T(v)=0
    \qquad\text{and}\qquad
    T(\alpha u)=\alpha T(u)=0.
  \]
  Thus $\ker T$ is a subspace. If $y_1=T(v_1)$ and $y_2=T(v_2)$ are in
  $\operatorname{im}T$, then
  \[
    y_1+y_2=T(v_1+v_2)
    \qquad\text{and}\qquad
    \alpha y_1=T(\alpha v_1),
  \]
  so the image is also a subspace.
\end{proof}

\begin{example}
  The continuous real-valued functions on $[0,1]$ form a subspace
  \[
    C([0,1])\subseteq\{f:[0,1]\to\bb{R}\}
  \]
  of the vector space of all real-valued functions on $[0,1]$.
\end{example}

\begin{example}
  Fix $b\in\bb{F}$ and let
  \[
    U_b=\left\{
      \begin{bmatrix}x_1\\x_2\\x_3\\x_4\end{bmatrix}\in\bb{F}^4:
      x_3=5x_4+b
    \right\}.
  \]
  Then $U_b$ is a subspace of $\bb{F}^4$ if and only if $b=0$. Indeed,
  the zero vector belongs to $U_b$ only when $b=0$, and when $b=0$ the
  defining equation $x_3=5x_4$ is preserved by addition and scalar
  multiplication.
\end{example}

\subsection{Sums of subspaces}

\begin{definition}[Sum of subspaces]
  If $V_1,\ldots,V_m$ are subspaces of $V$, their \emph{sum} is
  \[
    V_1+\cdots+V_m
    =\{v_1+\cdots+v_m:v_i\in V_i\text{ for }i=1,\ldots,m\}.
  \]
\end{definition}

\begin{example}
  Let
  \[
    U=\left\{\begin{bmatrix}x\\0\\0\end{bmatrix}:x\in\bb{F}\right\},
    \qquad
    W=\left\{\begin{bmatrix}0\\y\\0\end{bmatrix}:y\in\bb{F}\right\}.
  \]
  Then
  \[
    U+W
    =\left\{\begin{bmatrix}x\\y\\0\end{bmatrix}:x,y\in\bb{F}\right\}.
  \]
\end{example}

\begin{proposition}
  If $V_1,\ldots,V_m$ are subspaces of $V$, then
  $V_1+\cdots+V_m$ is the smallest subspace of $V$ containing each
  $V_i$.
\end{proposition}

\begin{proof}
  The sum contains $0=0+\cdots+0$. Adding two elements of the sum, or
  multiplying one by a scalar, can be done componentwise inside each
  $V_i$, so the sum is a subspace. For each $i$, any $v_i\in V_i$ can be
  written as
  \[
    v_i=0+\cdots+0+v_i+0+\cdots+0,
  \]
  and hence $V_i\subseteq V_1+\cdots+V_m$.

  Conversely, let $W$ be any subspace of $V$ containing all the $V_i$.
  If $v=v_1+\cdots+v_m$ with $v_i\in V_i$, then every $v_i$ belongs to
  $W$. Since $W$ is closed under addition, $v\in W$. Therefore
  $V_1+\cdots+V_m\subseteq W$.
\end{proof}

\subsection{Direct sums}

\begin{definition}[Direct sum]
  Let $V_1,\ldots,V_m$ be subspaces of $V$. Their sum is a
  \emph{direct sum} if every element of $V_1+\cdots+V_m$ has a unique
  representation
  \[
    v=v_1+\cdots+v_m,
    \qquad v_i\in V_i.
  \]
  In this case we write
  \[
    V_1\oplus\cdots\oplus V_m.
  \]
\end{definition}

\begin{proposition}[Zero-sum criterion]
  The sum $V_1+\cdots+V_m$ is direct if and only if the equation
  \[
    0=v_1+\cdots+v_m,
    \qquad v_i\in V_i,
  \]
  implies $v_1=\cdots=v_m=0$.
\end{proposition}

\begin{proof}
  If the sum is direct, then $0=0+\cdots+0$ is one representation of
  the zero vector, so uniqueness forces every other representation to
  have all components equal to zero.

  Conversely, suppose the zero vector has only the trivial
  representation. If
  \[
    v_1+\cdots+v_m=u_1+\cdots+u_m,
  \]
  where $v_i,u_i\in V_i$, then
  \[
    0=(v_1-u_1)+\cdots+(v_m-u_m).
  \]
  The hypothesis gives $v_i-u_i=0$ for each $i$, so $v_i=u_i$. Thus
  every representation is unique.
\end{proof}

\begin{proposition}
  If $U$ and $W$ are subspaces of $V$, then
  \[
    U+W=U\oplus W
    \quad\Longleftrightarrow\quad
    U\cap W=\{0\}.
  \]
\end{proposition}

\begin{proof}
  Suppose first that $U+W$ is a direct sum, and let $v\in U\cap W$.
  Then
  \[
    0=v+(-v)=0+0
  \]
  gives two representations with one term in $U$ and one term in $W$.
  Uniqueness implies $v=0$.

  Conversely, assume $U\cap W=\{0\}$. If
  \[
    u+w=u'+w',
    \qquad u,u'\in U,\quad w,w'\in W,
  \]
  then
  \[
    u-u'=w'-w\in U\cap W.
  \]
  Hence $u=u'$ and $w=w'$, so the representation is unique.
\end{proof}

For the coordinate-axis subspaces in the preceding example,
$U\cap W=\{0\}$, and hence their sum is $U\oplus W$.

\subsection{Product vector spaces}

\begin{definition}[Product vector space]
  Let $V_1,\ldots,V_m$ be vector spaces over $\bb{F}$. Their Cartesian
  product is
  \[
    V_1\times\cdots\times V_m
    =\{(v_1,\ldots,v_m):v_i\in V_i\}.
  \]
  Addition and scalar multiplication are defined coordinatewise:
  \begin{align*}
    (v_1,\ldots,v_m)+(u_1,\ldots,u_m)
      &=(v_1+u_1,\ldots,v_m+u_m), \\
    \alpha(v_1,\ldots,v_m)
      &=(\alpha v_1,\ldots,\alpha v_m).
  \end{align*}
\end{definition}

With these operations, $V_1\times\cdots\times V_m$ is a vector space
over $\bb{F}$; all vector-space axioms hold coordinatewise.

\begin{example}
  In $\mathcal P_5(\bb{R})\times\bb{R}^3$,
  \begin{align*}
    &\left(5-6t+4t^2,
      \begin{bmatrix}3\\8\\7\end{bmatrix}\right)
    +\left(t+9t^5,
      \begin{bmatrix}2\\2\\2\end{bmatrix}\right) \\
    &\qquad=
    \left(5-5t+4t^2+9t^5,
      \begin{bmatrix}5\\10\\9\end{bmatrix}\right).
  \end{align*}
\end{example}
